

\[
\boxed{\textbf{Dizemos que } |x| > 3 \textbf{ se:} \\
\\
\textbf{A.} \; x \in ] -\infty, -3[ \cup ]3, +\infty[ 
\quad\quad
\textbf{B.} \; x \in \mathbb{R} 
\quad\quad
\textbf{C.} \; x \in ] -3, 3[ 
\quad\quad
\textbf{D.} \; x \in ] -\infty, -3] \cup [3, +\infty[ 
\quad\quad
\textbf{E.} \; x \in ]3, +\infty[ }
\]





\[
\text{A desigualdade } |x| > 3 \text{ significa que o x está fora do intervalo } [-3, 3], \text{ ou seja, nos extremos da reta real.}
\]





\[
\text{Portanto, temos duas possibilidades: } x < -3 \text{ ou } x > 3.
\]





\[
\text{Em notação de conjuntos: } x \in ] -\infty, -3[ \cup ]3, +\infty[
\]





\[
\color{red}{\Huge \boxed{x \in ] -\infty, -3[ \cup ]3, +\infty[}}
\]


